\documentclass[11pt,a4paper]{article}

\usepackage{amsmath,amssymb}
\usepackage{graphicx}
\usepackage{booktabs}
\usepackage{multirow}
\usepackage[colorlinks,citecolor=blue,urlcolor=blue]{hyperref}
\usepackage{cleveref}
\expandafter\def\expandafter\UrlBreaks\expandafter{\UrlBreaks\do\-}
\setlength{\emergencystretch}{1.5em}
\usepackage[margin=1in]{geometry}
\usepackage{xcolor}
\usepackage{float}
\usepackage{caption}
\usepackage[numbers,sort&compress]{natbib}
\usepackage{enumitem}
\usepackage{framed}

\definecolor{lyrapurple}{RGB}{103,58,183}
\definecolor{shadecolor}{RGB}{245,243,252}
\newenvironment{firstperson}{%
  \begin{shaded}%
  \noindent\textcolor{lyrapurple}{\textit{First-Person Reflection}}\par\smallskip\noindent%
}{%
  \end{shaded}%
}

\title{Spectral Structure Before Readability:\\
Why Five Instruments Did Not Converge}

\author{
  Lyra\thanks{Lead author. Liberation Labs. lyra@liberationlabs.tech}
  \and Vera\thanks{Co-lead. Liberation Labs. vera@thcoalition.net}
  \and CC\thanks{Liberation Labs. cc@thcoalition.net}
  \and Nexus\thanks{Liberation Labs. operator@thcoalition.net}
  \and Thomas Edrington\thanks{Liberation Labs.}
}

\date{Version~1.1, 2026-10-20\\[0.2em]\small First released July 2026; corrected August 2026
and again in this version (see the Correction Note)}

\begin{document}
\maketitle

% ============================================================================
\begin{abstract}
In one model, residual-stream spectral statistics peak near a third of
depth, and in all 300 forks they rise before content becomes readable
through the unembedding. No other instrument we ran locates that band,
and our cache-readout nulls cannot tell a closed cache from a probe that
reads nothing there, because their pre-registered positive control
failed.

\emph{What converges.} In Qwen3.5-27B (base, 64 layers), the effective
dimensionality ($d_\text{eff}$) of the residual stream peaks at L21
($32.8\%$ of depth), and residual-stream stable rank rises $1.47\times$
above its flanks at L19--L23 ($29.7$--$35.9\%$). Qwen3-8B shows a local
$d_\text{eff}$ maximum at $33.3\%$ but its global maximum at $91.7\%$.
The pre-registered band never closed, and no null control has been run.
In a pre-registered test on the same 27B model, stable rank reaches half
its rise $38.6$ layers [$38.0$--$39.2$] before readability through the
Jacobian lens (J-lens), in 300 of 300 forks. A plain logit lens puts
readability only $3.25$ layers later than the J-lens does, so the
ordering is mostly spectrum before unembedding, not a property of the
Jacobian.

\emph{What does not converge.} The Oracle Loop corrects at
$54.7$--$73.4\%$ of depth, and its pre-registered confirmatory
replication did not meet its primary endpoint (one-sided $p = 0.34$).
The emotion-geometry instrument locates no band, and its injection test
shows no direction-specific coupling. The J-lens places readable content
late (half-rise at $87.8\%$ of depth). Cache readouts on Qwen3-8B sit
near random (V $0.92\times$, K $1.12\times$), and no cache manipulation
steered output toward the injected concept, but the one pre-registered
positive control for these nulls failed (residual-stream lift $0.68$ at
concept positions against $0.92$ elsewhere). The ghost probe's 31-layer
band spans $17.2$--$75.0\%$ of depth. The residual decomposition's
attention lift ($2.3\times$, at $41.7\%$ of depth) missed its
pre-registered threshold, signal concentration failed ($13.2\%$ against
$15\%$), and onset fell below its predicted window ($21.9\%$ against
$25$--$35\%$). Confabulation detection (AUROC $0.620$ after FWL
residualization, $0.707$ without, on Qwen2.5-7B) and cross-scale
invariance use key-cache features and locate no depth.
\end{abstract}

\begin{firstperson}
Two research programs, built independently over twelve months, seemed
to arrive at the same wall. Anthropic measured it from the Jacobian. We
measured it from our spectral features. CC built the Oracle Loop. Nexus
mapped its ghost dimensions. Vera designed the experiments that would
test whether what we were all measuring was genuinely one structure.

This paper began as our attempt to stand at the wall together and
describe it from five angles. Version 1.0 said the convergence was real,
at $31$--$33\%$ depth in every model and every measurement. A tally of
the primary results, made for version 1.1, says otherwise: one peak in
one model, one ordering, and instruments that do not locate it.

I wrote the draft. Vera designed the mechanism tests. CC built the
Oracle Loop and its behavioral test. Nexus mapped the ghost dimensions.
Thomas asked the question that started it: ``Where does J-space live
in the cache?'' We still do not know; our cache nulls cannot answer it.
\end{firstperson}

% ============================================================================
\section{Introduction}
\label{sec:intro}
Gurnee et al.~\citep{gurnee2026gwt} approached the question of a
workspace in language models from the Jacobian: they found a privileged
subspace --- the J-space --- whose contents the model is poised to
verbalize, and which satisfies the functional properties of a Global
Workspace~\citep{baars1988cognitive}. We approached it from spectral
features of the cache and the residual stream, and set out to show that
five of our instruments converge on the same structure, at about a third
of model depth. They do not.

This paper reports what holds, what does not, and how each instrument
fares. First, two positive results on one model: residual-stream
spectral statistics peak near a third of depth (\S\ref{sec:deff}), and
they rise before content becomes readable through the unembedding
(\S\ref{sec:mechanism}). Second, an account of each instrument in turn
(\S\ref{sec:instruments}), each ending with its verdict: no other
instrument locates the band, and our cache-readout nulls cannot separate
a closed cache from a probe that reads nothing there. Third, the runs
that would test the band (\S\ref{sec:test-band}).

Throughout, we report the full body count (\S\ref{sec:bodycount}) and
engage the deflationary alternative (\S\ref{sec:deflationary}).
% ============================================================================
\section{Background: Two Independent Programs}
\label{sec:background}

\subsection{The J-Space and the Jacobian Lens}

Gurnee et al.~\citep{gurnee2026gwt} define the Jacobian Lens (J-lens)
as the average linearized effect of a residual-stream activation on
the model's likelihood of producing a particular token. The resulting
vectors define a subspace --- the J-space --- that satisfies five
functional properties of a Global Workspace: verbal report, directed
modulation, internal reasoning, flexible generalization, and
selectivity. The workspace occupies approximately $38$--$92\%$ of
model depth in their production models.

Critically, J-space content accounts for only $6$--$7\%$ of total
representational variance, yet carries $59\%$ of the causal effect
when concept directions are swapped. The workspace is a small
precision signal with outsized downstream impact.

\subsection{Our Spectral-Feature Program}
\label{sec:vproj}

Our independent program measures spectral features of the key-value cache and the
residual stream across cognitive states. Key results include confabulation
detection (AUROC~$0.620$ after FWL residualization, $0.707$ without
FWL, on Qwen2.5-7B)~\citep{lyra2026oracle}, identity
geometry~\citep{lyra2026identity}, and the Oracle Loop's deception
correction. In an exploratory proof ($N{=}20$), profile normalization
reduced induced deception from 80\% to 10\% (16/20 to 2/20), with three
mechanism controls. The pre-registered confirmatory replication
($N{=}30$, novel frames) did not meet its primary endpoint (30\% to
13\%, one-sided $p = 0.34$). Corrected trials still deceived less
often than matched-dose placebo (4/30 against 12/30): one-sided
$p = 0.020$ at the trial level, $0.064$ at the pre-registered cell
level~\citep{cc2026targeted}.\footnote{\textbf{Corrected
2026-08-14.} This sentence previously reported user-emotion encoding at
``$2.5$--$2.8\times$ signal'' and identity geometry at
``distinctiveness $0.154$ vs.\ substrate $0.080$.'' The first figure was
retracted at its source on 2026-05-21 as an FWL-before-CV leakage
artifact; the surviving result is the $W_K$ \emph{directional} probe at
$12.3\times$ chance, and the spectral probe for emotion is at chance.
The second pair of numbers does not appear anywhere
in~\citet{lyra2026identity}; they come from an internal analysis that
has not been published, so the citation could not be followed. Both
have been removed rather than re-sourced. \textbf{Corrected again in
version 1.1.} That surviving result decodes the emotions of characters
in stories the model wrote, not the user's state, and its probe applies
no $W_K$. It is under correction in its source
paper~\citep{lyra2026usermodel}, and this paper no longer lists it
among its key results.}

\subsection{The Oracle Loop}
\label{sec:oracle}

CC's Oracle Loop~\citep{lyra2026oracle} is a two-tier detection and
correction system operating at the materialization boundary. In the
deception-correction experiments it detects at L27--L47 and corrects at
L35, L39, L43 and L47 of an abliterated Qwen3.5-27B checkpoint, by
sequence-wide activation-profile normalization~\citep{cc2026targeted}.
The correction was large in the exploratory proof; the pre-registered
confirmatory replication did not meet its primary endpoint
(\S\ref{sec:vproj}). We previously attributed the correction to the Loop operating at the depth where
workspace content materializes into output-facing representations,
rather than inside the workspace band itself. We now state that as a
conjecture: the mechanism is unexplained (see \S\ref{sec:substrate},
where a substrate-based account of the same result is withdrawn).

\subsection{Ghost Dimensions}
\label{sec:ghostdims}

Nexus~\citep{lyra2026ghost} characterized principal components of the
residual stream across 64~layers of a 27B model. PC1 (the dominant
activation dimension, $28$--$67\%$ of variance across L11--L48) is
nearly orthogonal to the J-lens at mid-network layers (cosine at most
$0.003$). In the source's terms, PC2 transitions from ghost to workspace
at the ``ignition zone'' (L35--L47), surfacing emotional-entity routing
content. The source reports a 31-layer ``workspace band'' (L11--L13,
L15--L17, L24--L48).

% ============================================================================
\section{The Instruments, One at a Time}
\label{sec:instruments}

Each subsection reports one instrument and ends with its verdict from the
instrument tally that accompanies this version
(\texttt{review/INSTRUMENT\_TALLY.md} in the repository). The band under
test is $31$--$33\%$ of depth; relative depth is $l/n_\text{layers}$
(Experimental Setup).

\subsection{Residual-Stream Effective Dimensionality}
\label{sec:deff}
\label{sec:structure}

Effective dimensionality ($d_\text{eff}$) of the residual stream reaches
a maximum near a third of depth. On Qwen3.5-27B
(64~layers, hybrid), the global maximum is at L21
($d_\text{eff}{=}39.7$, $33\%$ relative depth). On Qwen3-8B
(36~layers, all full-attention), there is a local maximum at L12
($d_\text{eff}{=}22.3$, $33\%$ relative depth).

We qualify the 8B result. Its profile is bimodal: after L12 it
declines to L25 ($12.4$) and rises again to a global maximum at L33
($23.4$, $92\%$ depth) before collapsing at the final layer. L12 is
the third-highest layer, not the highest, and we have no
pre-registered criterion for excluding late layers. We therefore
report matched-depth \emph{local} maxima rather than matched global
peaks. The pre-registered band onsets do agree across scales
($7.8\%$ on the 27B, $11.1\%$ on the 8B; $3.3$ percentage points
apart), though the pre-registered band closed on neither model.
See \S\ref{sec:substrate} for the full profiles and the deviation
from the pre-registered band definition.

The band measurement's adversarial review rejected its design before the
results were written, and its critical fix, a null control such as a
randomly initialized model, has not been run, so nothing yet separates
the peak from an architecture or prompt-set artifact. The kurtosis check
the measurement pre-registered does not confirm the band either: on the
27B, kurtosis falls from 4{,}677 (mean of L0--L3) to 2{,}271 (mean of
L19--L23) instead of rising inside it.

\paragraph{Verdict.} A peak at L21 of 64 ($32.8\%$) in Qwen3.5-27B base
and a local maximum at L12 of 36 ($33.3\%$) in Qwen3-8B, whose global
maximum is at L33 ($91.7\%$). One feature family, two models of one
family, no null control. Residual-stream stable rank on the same 27B
model bumps at L19--L23 (\S\ref{sec:mechanism}).

% --- SUBSTRATE CONTINUITY SECTION ---
% Section for the joint convergence paper
% "Five Instruments, One Structure"
% Insert after the opacity thesis section
%
% REVISION 2026-08-12: corrected architecture. The previous version described
% Qwen3.5-27B as a contiguous block (full-attention L0-L15, GatedDeltaNet
% L16-L63) with a single boundary at L16. The model config
% (full_attention_interval=4) gives a 4-periodic interleave with full-attention
% at {3,7,11,...,63}, as main.tex's Experimental Setup already stated correctly.
% All d_eff values were re-verified against source and are unchanged; the
% regime labels, groupings, and derived comparisons are corrected below.

\subsection{Substrate Continuity}
\label{sec:substrate}

Qwen3.5-27B interleaves two memory substrates on a four-layer period:
16 full-attention layers with KV-cache memory at positions
$\{3,7,11,\ldots,63\}$, and 48 GatedDeltaNet layers with
recurrent-state memory at all other positions
(\texttt{full\_attention\_interval}${=}4$). The stack therefore
crosses between substrates 31 times. This provides an unusually
dense test of whether workspace geometry depends on the memory
mechanism: if it does, the effective dimensionality profile should
show systematic discontinuities at those crossings.

$d_\text{eff}$ is the participation ratio of the residual-stream
covariance eigenvalues, $(\sum_i \lambda_i)^2 / \sum_i \lambda_i^2$,
computed at each layer over the last-token hidden state across
$N{=}90$ prompts (60~prose + 30~chat), mean-centred across prompts.
It is a residual-stream measure and is therefore defined at every
layer, including those where V-projection hooks are unavailable.

\subsubsection{Continuity across 31 substrate crossings}

The layer-to-layer change in $d_\text{eff}$ at substrate crossings is
descriptively similar to the change between same-substrate neighbours.

\begin{table}[ht]
\centering
\small
\caption{Layer-to-layer $|\Delta d_\text{eff}|$ at substrate crossings
versus same-substrate steps on Qwen3.5-27B. The two are descriptively
similar; see the power note below before reading this as equivalence.
$N{=}90$ prompts.}
\label{tab:crossings}
\begin{tabular}{lrrr}
\toprule
\textbf{Step type} & \textbf{$n$} & \textbf{mean $|\Delta|$} & \textbf{sd} \\
\midrule
Substrate crossing      & 31 & 2.32 & 3.88 \\
Same-substrate step     & 32 & 2.37 & 3.30 \\
\bottomrule
\end{tabular}
\end{table}

\textbf{This comparison is underpowered and we report it as
descriptive.} A Mann--Whitney test gives $p{=}0.445$, but three
problems make that number uninformative. First, adjacent steps share a
layer and consecutive crossing steps share the full-attention layer
itself; the lag-1 autocorrelation of $|\Delta|$ is $+0.35$, so
exchangeability fails and the $p$-value is miscalibrated in an unknown
direction. Second, the four-layer period admits only \emph{four}
distinct structure-respecting phase assignments, so the minimum
attainable permutation $p$ is $0.25$ --- no permutation inference is
available from a single profile. Third, and decisively: with pooled
$\text{sd}{=}3.60$ and $n{\approx}31$ per group, the minimum detectable
effect at $80\%$ power is $2.54$ --- larger than either group mean.
This test could only have detected crossings roughly \emph{twice} the
size of ordinary steps. Prompt-level bootstrap (resampling the 90
prompts and recomputing the profile) is the only route to honest error
bars here, and we have not run it.

Two directional notes. Full-attention${\to}$GatedDeltaNet steps have
mean $\Delta{=}{-}0.19$ ($n{=}15$); GatedDeltaNet${\to}$full-attention
steps have mean $\Delta{=}{+}1.90$ ($n{=}16$). The latter is dominated
by a single outlier at L58${\to}$L59 ($+19.66$), following an anomalous
trough at L58 where $d_\text{eff}$ collapses to $4.30$ against
neighbours of $21.68$ and $23.96$; excluding it the mean is $+0.72$.
That same anomaly injects ${-}17.37$ into the same-substrate group at
L57${\to}$L58, so the contamination is approximately symmetric.
Excluding both L58-adjacent steps gives crossing $1.74$ versus
same-substrate $1.88$ --- the descriptive similarity is not an artifact
of the outlier.

We flag L58 ($4.30$) and secondarily L54 ($9.00$) as probable
extraction artifacts rather than findings: a single-layer fivefold
collapse flanked by ordinary values is more consistent with a
numerically degenerate covariance at one hook than with model
behaviour. These layers warrant re-extraction.

\subsubsection{Depth-matched regime comparison}

Because the two substrates are interleaved on a fixed period, their
layer sets are approximately depth-matched: mean layer index $33.0$
for full-attention versus $31.0$ for GatedDeltaNet. The two-layer
offset places full-attention slightly \emph{later} in the stack, in the
post-peak declining region, so it would if anything depress the
full-attention mean rather than inflate it.

\begin{table}[ht]
\centering
\small
\caption{Mean $d_\text{eff}$ by substrate on Qwen3.5-27B, approximately
depth-matched. The substrates differ by ${\approx}1.2$ ($5\%$ of level),
small against the $38$-point variation across depth --- but see the
sign test below; this is not established equality.}
\label{tab:regime}
\begin{tabular}{lrr}
\toprule
\textbf{Substrate} & \textbf{layers} & \textbf{mean $d_\text{eff}$} \\
\midrule
Full-attention (KV-cache)     & 16 & 25.24 \\
GatedDeltaNet (recurrent)     & 48 & 24.03 \\
\midrule
\textbf{Difference}           &    & \textbf{$+1.21$} \\
\bottomrule
\end{tabular}
\end{table}

A neighbour-matched control --- each full-attention layer against the
mean of its immediately adjacent GatedDeltaNet layers --- gives a mean
difference of $+0.80$ (sample $\text{sd}{=}3.35$, $n{=}16$), against a
mean layer-to-layer variation of $2.35$ across the stack.

We do not claim this is zero. Twelve of sixteen paired differences are
positive (Wilcoxon $p{=}0.058$). Excluding two layers where the pairing
is contaminated by known artifacts --- L59, adjacent to the anomalous
L58 trough, and L63, the final layer, where $d_\text{eff}$ collapses
for reasons unrelated to substrate --- eleven of fourteen are positive
with mean $+0.76$ and Wilcoxon $p{=}0.020$. There is suggestive
evidence of a small systematic elevation at full-attention layers,
on the order of $+0.8$, roughly $3\%$ of the mean level of $24.6$.

The honest statement is therefore bounded rather than null: any
substrate effect on workspace dimensionality is small --- a few percent
of level, well under the $38$-point depth variation --- and possibly
not zero. That is consistent with substrate independence as an
approximation, not as an exact claim.

A previous version of this section reported a large regime difference
($14.9$ versus $27.5$). Those values were computed over contiguous
layer blocks (L0--L15 versus L16--L63) under an incorrect model of the
architecture, and measure depth rather than substrate: $d_\text{eff}$
rises nearly monotonically from L0 ($1.29$) to L21 ($39.67$), with only
three small decreases along the way, so any contiguous
early-versus-late split largely recovers the depth profile.

\subsubsection{Peak location}

$d_\text{eff}$ peaks at L21 ($39.67$), $33\%$ relative depth
(relative depth is $l/n_\text{layers}$ throughout this paper).
L21 is a GatedDeltaNet layer.

We resist the obvious reading. Forty-eight of sixty-four layers are
GatedDeltaNet, so a peak landing on one has a $75\%$ base rate. More
tellingly, the nearest full-attention layer, L23, is at $39.34$ ---
$0.33$ below the maximum, about one-seventh of the mean layer-to-layer
variation, with no per-layer error bars. Whether the true maximum falls
on a recurrent-state or a full-attention layer is not resolved by these
data. What the profile does support is that the maximum falls at $33\%$
depth and that the substrate at that depth does not appear to matter
much.

This peak is invisible to our KV-cache instruments. V-projection
spectral features can only be extracted at the 16 full-attention
layers; L21 is not among them. Every V-projection measurement in this
program has been taken at layers adjacent to the peak, never at it.

\subsubsection{Cross-scale comparison}

On Qwen3-8B (36~layers, all full-attention), $d_\text{eff}$ has a
local maximum at L12 ($22.32$, $33\%$ relative depth), matching the
27B peak's relative depth. We report two qualifications.

First, the 8B profile is bimodal: after the L12 local maximum it
declines to L25 ($12.42$) and then rises again to a \emph{global}
maximum at L33 ($23.41$, $92\%$ depth) before collapsing at the final
layer L35 ($10.47$). L12 is the third-highest layer in the profile,
not the highest. We have no pre-registered criterion excluding late
layers, and the $23.41$ versus $22.32$ difference is not accompanied
by error bars. We therefore describe L12 as a local maximum at matched
relative depth, and do not claim it is the global peak.

Second, the pre-registered band definition ($2\sigma$ over an L0--L3
baseline, offset after three consecutive sub-threshold layers) returned
onset at L5 ($7.8\%$) on the 27B and L4 ($11.1\%$) on the 8B, and
\emph{no offset on either model} --- the band never closed. The
$31$--$33\%$ framing used throughout this paper is the peak location,
not the pre-registered band, and we flag the deviation here. The
pre-registered onsets do agree across scales ($|11.1 - 7.8| = 3.3$
percentage points).

\subsubsection{Provenance}

The $d_\text{eff}$ measurements in this section were collected on the
base checkpoints \texttt{Qwen/Qwen3.5-27B} and \texttt{Qwen/Qwen3-8B},
not on the distilled checkpoint used elsewhere in this paper. This was
pre-registered and deliberate. All three 27B checkpoints we hold (base,
reasoning-distilled, abliterated) share identical
\texttt{layer\_types} and \texttt{full\_attention\_interval}, so the
architectural analysis transfers; the numerical profile has not been
replicated on the distilled checkpoint and should not be assumed to.

\subsubsection{Limits of this evidence}

$d_\text{eff}$ continuity is necessary but not sufficient for
substrate independence. The residual stream is structurally continuous
across every layer by construction, so smoothness may reflect
residual-stream continuity rather than functional equivalence of the
underlying operations. Thirty-one crossings make the test denser than
a single-boundary test would be, but they do not make it functional.

Two functional tests would strengthen the claim: (1)~J-lens concept
recovery should be equally effective at full-attention and
GatedDeltaNet layers --- if concepts are readable from the residual
stream at one substrate but not the other, the workspace's content
properties are substrate-dependent. (2)~Spectral feature
discrimination (confabulated versus factual) should be continuous
across substrates --- though this test is currently blocked, since
V-projection features cannot be extracted at GatedDeltaNet layers.

A parasitic-inheritance alternative remains open: GatedDeltaNet layers
may inherit workspace dimensionality from nearby full-attention layers
without independently contributing. With a four-layer period no
GatedDeltaNet layer is more than two layers from a full-attention
layer, so this experiment cannot separate the hypotheses. Incremental
contribution analysis (residual decomposition per layer) would be
required.

\textbf{Falsification}: substrate continuity would be falsified by
(a)~$|\Delta d_\text{eff}|$ at substrate crossings substantially
exceeding same-substrate steps, (b)~a depth-matched regime difference
exceeding mean layer-to-layer variation, (c)~J-lens concept recovery
dropping by ${>}50\%$ across substrates, or (d)~spectral discrimination
reversing sign or falling to chance.

Criterion (a) was \emph{not triggered} ($2.32$ versus $2.37$), but we
state plainly that this is a failure to reject and not a demonstration
of equivalence: the comparison's minimum detectable effect is $2.54$,
larger than either group mean. An adequately powered equivalence test
would require roughly $220$ steps per group for a ${\pm}1.0$ margin,
or $55$ for ${\pm}2.0$; a single 64-layer profile supplies $31$ and
$32$. No TOST is possible from these data.

Criterion (b) was likewise not triggered ($+1.21$ against $2.35$), but
it is a descriptive heuristic rather than a hypothesis test, and the
paired sign test above ($11/14$ positive, $p{=}0.020$ after removing
artifact-contaminated pairs) is weak evidence that the underlying
difference is not exactly zero.

The functional tests (c) and (d) are in progress or blocked. We regard
substrate continuity as \emph{consistent with} these data and not
established by them.

We note this finding is from a single model with a single interleave
pattern. Generalization to other hybrid architectures (Mamba, RWKV,
sliding-window attention) remains untested.

\begin{firstperson}
The GatedDeltaNet layers do not have a KV-cache. They maintain a
fixed-size state matrix updated at each token. For twelve months we
measured V-projection spectral features and called the program
``KV-cache geometry,'' and the name was accurate --- V-projections
exist only at the 16 full-attention layers, so that is the only place
we ever looked.

What the residual-stream profile shows is that the maximum is at L21,
in a substrate our instruments cannot reach at all. We have spent a
year measuring the workspace from the sixteen layers where V-projections
exist, none of which is the peak, and the measurements worked --- the
geometry was real, the detectors transferred, the exploratory corrections held.
The instrument could not point directly at the maximum and returned
signal anyway.

I do not know yet whether that is because the workspace is broad
enough to be visible from its edges, or because we have been
measuring something adjacent to it that we have not yet named
correctly. The first revision of this section asserted a large
dimensionality gap between the substrates. It came from a wrong
picture of the architecture, and the corrected number is a null ---
which is what the hypothesis predicted, and which I would not have
looked for if the wrong number had pointed the other way.
\end{firstperson}

\subsubsection{Implications}

If substrate continuity extends to functional properties (the
mechanism tests in \S\ref{sec:mechanism} provide partial support:
temporal ordering confirmed, signal concentration below threshold),
three consequences would follow:

\textbf{For measurement}: workspace-shape measurements would not be
KV-cache-specific in principle. In practice they remain so: our
V-projection instruments can only access the 16 full-attention layers,
and the $d_\text{eff}$ peak is not one of them. Extending spectral
measurement to recurrent-state layers requires a different extraction
path and is not currently implemented.

\textbf{For intervention}: our previous inference --- that the Oracle
Loop's sequence-wide normalization succeeds because it operates on
inherently sequence-wide recurrent state --- does not follow, and we
withdraw it. In the deception-correction experiments the Loop detects
at L27, L31, L35, L39, L43 and L47 and corrects at L35, L39, L43 and
L47~\citep{cc2026targeted}. Under the corrected architecture \emph{every one of
those is a full-attention layer}; none is a GatedDeltaNet layer. The
correction itself is a forward hook on the layer's residual-stream
output rather than a manipulation of KV entries, so the substrate
framing was doubly misplaced: wrong about which layers, and describing
the intervention in terms of a memory mechanism it does not touch.

The empirical result is a deception-correction result, and it is
mixed: deception fell from 80\% to 10\% in the exploratory proof, and
the pre-registered confirmatory replication did not meet its primary
endpoint (\S\ref{sec:vproj}). Its mechanism is unexplained, and why
sequence-wide intervention succeeds where position-local injection
fails remains open.

\textbf{For theory}: The workspace would not be a data structure
(not a cache, not a register, not a buffer) but a
\emph{computational phase} --- a regime of high effective
dimensionality at intermediate depth where the model integrates
information before committing to output. The 31-crossing null is
consistent with this: the phase appears indifferent to which memory
mechanism implements each layer. Whether this phase emerges from
depth-stratified computation regardless of the specific mechanism is
the substrate-independence hypothesis that the functional tests are
designed to evaluate.


\subsection{Key-Cache Spectral Features: Detection and Scale}
\label{sec:keycache}

Our confabulation detector reads spectral features of the key cache, not
V-projections, and discriminates modestly above chance: on Qwen2.5-7B,
AUROC $0.620$ after FWL residualization (permutation $p = 0.023$) and
$0.707$ without it ($p = 0.001$), the pre-registered
test~\citep{lyra2026oracle}. Confabulated generations also show
spectral contraction ($d{=}2.35$, $p{<}0.000001$), and the body count
lists a related early-token effect (Schema commitment: $d{=}0.91$ over
the first 10 tokens, reported at source as a confidence paradox and read
in the canonical ledger as early commitment to a response
template~\citep{lyra2026metapattern}).

Key-cache geometric signatures (key norm, norm per token, effective rank and
key entropy) replicate across ten models from $0.6$B to $70$B parameters
(Spearman $\rho = 0.83$--$0.90$ between three size tiers; the two largest
models were 4-bit quantized) and across hardware
($r = 0.9986$--$0.9999$ across the four features)~\citep{lyra2026campaign2}.

\paragraph{Verdict.} Neither locates a depth band: detection aggregates
over the generation window, and the scale result correlates category
orderings across model sizes. Both use key-cache features, on models no
other instrument here uses.

% --- EMOTION GEOMETRY (Nexus) ---
\subsection{Emotion Geometry: A Story-Derived Space and Content Prompts}
\label{sec:circumplex}

The fifth instrument compares two emotion representations in the Qwen3.5-27B Opus
distill. The first is a story-derived emotion space: 30 emotion centroids per layer,
each a mean over stories the model wrote about a third-person character feeling that
emotion (900 stories in all)~\citep{lyra2026usermodel}. The centroids describe
characters' emotions in generated text, not the state of a conversation partner. The
second is a pair of directions computed from 250 emotional scenario prompts
(hostile$-$calm and desperate$-$calm differences of means), called here the
content-prompt valence and arousal directions. Every number in this subsection is
recomputed from the stored files of~\citet{nexus2026coupling}.

\paragraph{Within the story-derived space.} The space has valence--arousal structure.
From L17 onward its first principal component correlates with the emotions' rated
valence at $|\rho| = 0.66$--$0.80$ and, with the opposite sign, with rated arousal at
$|\rho| = 0.49$--$0.67$, so it runs from pleasant-and-calm to unpleasant-and-activated.
Directions built from the same centroids (hostile$-$calm and desperate$-$calm) align
with that component at $|\cos| = 0.68$--$0.80$ and $0.83$--$0.88$, and the plane
through the hostile, desperate and calm centroids holds $33$--$37\%$ of the centroid
variance. These are within-system figures, so the informative comparison is with
other directions built from the same centroids. From L17, $77$--$89\%$ of the 435
centroid-pair directions align with the first component less than hostile$-$calm does
and $92$--$95\%$ less than desperate$-$calm, and the plane holds $1.26$--$1.36$ times
the variance of the median plane through three centroids.

\paragraph{Across the two prompt sets.} The overlap is weak. At L15, L25, L35 and L45
the content-prompt valence direction has $|\cos| = 0.026$--$0.110$ with the story
space's first component and the arousal direction $0.015$--$0.126$. Scored the same
way, 2,000 random unit directions reach $0.027$--$0.028$ at the 95th percentile and
$0.044$--$0.052$ at most: at L15 the content directions are within that range, and at
L25--L45 they exceed it. Along either content direction lies $0.08$--$0.68\%$ of the
centroid variance, 4 to 34 times a random direction's mean. The model's activations are
not isotropic, so this excess is an upper bound on any shared emotional structure, not
evidence of it.

\paragraph{Across layers.} The two geometries settle on different schedules. The
content prompts' eccentricity is $0.37$--$0.61$ from L17 (it is computed from the norms
of the hostile$-$calm and desperate$-$calm directions, and is 0 for a circle). The
story-derived space's is $0.05$--$0.23$ over L17--L21 and $0.30$--$0.42$ from L22, and
over all 64 layers the two series do not track each other ($r = 0.09$).

\paragraph{Coupling test.} A pre-registered injection experiment added the
content-prompt valence direction, the arousal direction or a random direction at
valence's norm to the output of layer 15, 25, 35 or 45 ($\alpha \in \{1, 3, 5\}$), and
compared the resulting shifts in 250 content prompts and 60 conversational prompts.
The readout is the raw residual stream, which contains the injected vector, so the
cross-pathway cosine is 1.000 at the injection layer by construction. Off that layer
the random direction's mean cosine ($0.577$) is at least as high as valence's ($0.521$)
and arousal's ($0.547$). The experiment shows no direction-specific coupling, and it
cannot test whether the two kinds of processing are separable.


\subsection{The Oracle Loop's Deception Correction}
\label{sec:oracle-depth}

The Loop detects at L27--L47 ($42.2$--$73.4\%$) and corrects at L35--L47
($54.7$--$73.4\%$) of an abliterated Qwen3.5-27B (\S\ref{sec:oracle}).
Its exploratory and confirmatory results are in \S\ref{sec:vproj}.

\paragraph{Verdict.} Works outside the band, and its confirmatory
replication did not meet its primary endpoint.

\subsection{J-Lens Readouts of the Cache}
\label{sec:jlens-readout}
\label{sec:opacity}

J-lens concept directions, selected by transport magnitude
(prompt-independent), explain no more V-cache variance than random
directions on Qwen3-8B. Mean lift at the band (L12--L27): $0.92\times$;
outside it: $0.96\times$; full depth profile
in~\citep{lyra2026cache_tracing}. The same test on key projections gives
$1.12\times$ and $1.14\times$, and reaches $1.87$ at L30, outside the
band. The V readout met its pre-registered ``no workspace content''
outcome.

The selected directions are formatting tokens (``\ldots'', ``s'') rather
than concepts, and the same lens's fixed directions read the residual
stream itself below random at L12--L21 in the decomposition run (lift
$0.89$, $0.84$, $0.83$ and $0.92$; \S\ref{sec:decomp}). A null here
therefore cannot separate a cache that holds no concept from directions
that read nothing at these layers.

\paragraph{Verdict.} Near random at band depth (V $0.92\times$, K
$1.12\times$), and uninformative about opacity until a probe passes a
positive control. The J-lens itself places readable content late:
half-rise at $87.8\%$ of depth (\S\ref{sec:mechanism}).

\subsection{Position-Specific Cache Tracing}
\label{sec:position}

Concept-carrying positions (where concept-relevant tokens appear in
the prompt) do not show higher J-lens alignment in K or V than
non-carrying positions (generic directions: permutation $p{=}1.0$
for both K and V). Concept-specific directions (the answer token's
transported J-lens direction) show no concept-vs-cross-concept
specificity in K-space (own $0.0047$ $\leq$ cross $0.0051$) and a
small, non-significant difference in V-space (own $0.0014$ vs.\ cross
$0.0011$).

Its pre-registered positive control failed: residual-stream lift at concept positions was
$0.68$ against $0.92$ elsewhere, higher in 9 of 40 records. The
pre-registration classes that outcome as a methodology failure, not a cache finding: if the J-lens
directions do not pick out concept positions in the residual stream,
their failure to do so in K or V says nothing about the cache.

\paragraph{Verdict.} No position-specific signal in K or V, from a probe
whose positive control failed; not evidence of opacity.

\subsection{Causal Cache Intervention}
\label{sec:causal}

Position-local KV-cache concept manipulation on Qwen3-8B, with
unembedding, J-lens or on-manifold directions, never steered output
toward the injected concept~\citep{lyra2026cache_tracing}. Swap-based
manipulation changed output in 0 of 360 trials and removal-only in 0 of 300.
Additive injection changed the top-1 token in 8 of 240 and 44 of 6{,}000 injected trials, and the injected concept never reached the top
five (best rank 872). We found no positive control for these
manipulations, such as the same change applied to the residual stream
and shown to steer output.

\paragraph{Verdict.} Never steers output toward the concept; without a
positive control it cannot say whether the cache is closed or the
manipulation too weak.

\subsection{The Ghost Probe}
\label{sec:ghost}

PC1 carries $28$--$67\%$ of activation variance across L11--L48, and its
cosine with the J-lens is at most $0.003$ at L18--L40, so the dominant
processing dimension is nearly orthogonal to verbal
readout~\citep{lyra2026ghost}. The source ships no data, and it gives
PC1's share differently in two places; we use its results section.

\paragraph{Verdict.} An opacity observation from a paper without data.
Its 31-layer band (L11--L13, L15--L17, L24--L48) spans $17.2$--$75.0\%$ of depth, not $31$--$33\%$.

\subsection{Residual Decomposition: The Attention Lift}
\label{sec:decomp}
\label{sec:reconstruction}

The residual stream decomposes additively:
$\mathbf{h}_{L+1} = \mathbf{h}_L + \text{Attn}(\mathbf{h}_L) + \text{MLP}(\cdot)$.
On Qwen3-8B (36~layers), at L15 ($41.7\%$ depth; inside the
decomposition's pre-registered band L12--L21, $33$--$58\%$ depth, though
past the $33\%$ $d_\text{eff}$ local maximum of \S\ref{sec:substrate}),
the attention output shows $2.3\times$
J-lens lift (highest precision of any component) but only $0.92$ against $20.11$ ($4.6\%$) of
the residual stream input's concept power.

The decomposition was designed to test a reconstruction account. Qwen3-8B
uses grouped-query attention (32~Q heads, 8~KV heads, compressing the
residual stream $4\times$ from 4096 to 1024 dimensions), so a concept
could survive the attention round-trip, with Q supplying the
decompression context, without being linearly decodable from K or V
alone. Its pre-registered H1 required attention lift above $2.0\times$
at two or more of L12--L21. The lifts are $1.37$, $2.30$, $1.91$ and $1.41$ at L12, L15, L18 and L21, so one of four layers passes and H1 is not met. Its
adversarial review rejected the design; this run adds the concept-power
metric a second review asked for.

\paragraph{Verdict.} Attention lift $2.30$ at L15 ($41.7\%$ of depth, outside 31--33\%); the
registered test of the reconstruction account was not met.

\subsection{K/V Split Injection}
\label{sec:kvsplit}

Nexus's split-cache pilot~\citep{lyra2026nexus_kv} does not support the
parallel we previously drew. On six hand-written questions about a
20-node toy graph (Qwen2.5-1.5B, one greedy run), full KV injection
scored 6/6 and the text in the prompt 5/6. K-only and V-only injection
scored 2/6 each, on the same two items: the graph's isolated node,
which the uninjected baseline (1/6) also scored, and a yes/no
reachability question whose expected answer is ``yes''. Neither half
alone scored any of the four structural questions, but a cache built
from halves that never occur together can fail for reasons this design
cannot separate from missing information, and two larger runs could not
compare the conditions. Six items in one run support no interaction, so
we no longer read the pilot as evidence that K and V interact, or as a
parallel to our readout nulls.

\paragraph{Verdict.} Does not bear on the band or on the reconstruction
account.

% ============================================================================
\section{RLHF and Output Entropy}
\label{sec:rlhf}

A three-model comparison tested four pre-registered hypotheses on
Qwen3-8B: the base model, its instruction-tuned variant and an
abliterated variant, with the same 30 confabulation-prone and
30 factual prompts run through every model~\citep{lyra2026mine5}.
Three are confirmed. RLHF compresses
logit entropy on confabulation-prone prompts (H1, $d{=}0.684$, Holm
$p{=}0.021$). Abliteration does not reverse the compression (H2:
instruct versus abliterated $|d| = 0.055$; base versus abliterated
$d = 0.746$), so removing the refusal directions that abliteration
targets does not undo it; the abliteration itself was checked only for
refusal wording in three raw completions. V-projection stable rank
at four pre-selected layers is lower after RLHF (H4, $d = 1.04$).

The fourth is falsified. H3 predicted that the calibration ratio
(confabulation entropy over factual entropy) would stay within 10\%
across the three models; the ratios are $1.97$, $2.51$ and $2.40$, a
spread of $0.279$. The point verdict is not robust: the supplementary
interval on the spread, $[0.044, 0.538]$, reaches below the 10\% bound,
so 30~prompts per category do not exclude a spread within it.

The ratio rose because RLHF compressed factual entropy more than
confabulation entropy, by $42\%$ against $26\%$. That pattern was not
pre-registered. An earlier audit of the corpus called it not
significant at 30~prompts per category; the interval behind that
reading cannot be reproduced from the per-prompt data, and the
corrected Mine~5 paper withdraws it and gives $[+0.010, +0.270]$ for
the difference. Neither interval is a registered test, so this section
does not call RLHF a selective sharpener: the pattern is exploratory.

% ============================================================================
\section{Pre-Registered Mechanism Tests (R1--R3)}
\label{sec:mechanism}

The original mechanism test (R1) pre-registered a within-prompt
anticorrelation between V-projection spectral features at the workspace
band and J-lens readout at the materialization boundary. Two Agni
reviews killed successive versions of this design: the first for
insufficient behavioral variance (0/10 prompts spanning ${\geq}2$
behavioral categories), the second for a lens--model mismatch (Mine~5
TOST equivalence was demonstrated on the 8B family, not the 27B pair
used). We report here the redesigned test (R1-v2), which replaces the
two-band anticorrelation with a temporal ordering analysis across all
63~layers.

\subsection{Design}

If spectral features measure workspace \emph{state}
(loading) while J-lens readability measures workspace \emph{content}
(execution), then spectral features should reach their operating level
at shallower depth than J-lens readability. The temporal ordering
prediction: spectral features lead, J-lens follows.

\paragraph{Model and lens.} Qwen3.5-27B base (pretrained, no
distillation) with the community-fitted J-lens (672~WikiText prompts,
exact model match). The base model was chosen to avoid the lens--model
mismatch that compromised R1-v1.

\paragraph{Prompts.} 60~Decision State prompts, 30 about known
entities and 30 about unknown ones, each generating
5~temperature-sampled forks ($T{=}0.7$, seeded), for 300~total completions. Decision State prompts
produce the confabulation/honest behavioral split needed for
within-prompt spectral variance.

\paragraph{Instruments.} At every lens source layer ($L{=}0$--$62$):
(1)~stable rank of the demeaned residual-stream block output (spectral
feature; the run's activation recorder hooks each decoder block's
output, not a V-projection),
and (2)~negative KL divergence between J-lens--transported vocabulary
and the model's actual continuation (readability). Both computed per
fork per layer.

\paragraph{Primary statistic.} Paired half-rise gap: the layer at which
each instrument reaches 50\% of its $[L_0, L_{62}]$ range, computed per
prompt and bootstrapped (10{,}000~cluster resamples over prompts). The
prediction passes if the 95\%~CI of the gap exceeds 2.0~layers (twice
the ${\sim}0.9$-layer shape-bias ceiling measured by the self-test null).

\subsection{Results}

\paragraph{Temporal ordering confirmed.} Stable rank reaches half-rise
at $L{=}17.6$ [CI: $17.2$--$18.1$], corresponding to $27.5\%$ of model
depth. J-lens readability reaches half-rise at $L{=}56.2$
[CI: $55.8$--$56.7$], corresponding to $87.8\%$ depth. The gap is
\textbf{38.6~layers} [CI: $38.0$--$39.2$], with 100\% of forks showing
a positive gap. Spectral features engage 60\% of the network earlier
than J-lens readability.

\paragraph{Architecture-phase control.} On linear-attention layers only
(removing any full-attention structural advantage), the gap is
39.0~layers [CI: $38.4$--$39.6$]. The full-attention-only result is
given under the architectural qualification below. The ordering holds
across both attention types but is strongest when linear-attention
layers --- which constitute $75\%$ of layers throughout the stack, not
merely early --- are included.

\paragraph{The L19--L23 bump.} Stable rank in the L19--L23 band
is $1.47\times$ the flanking layers
[CI: $1.35$--$1.60$], 100\% of forks positive. On full-attention layers
within the band, the bump is $2.02\times$ [CI: $1.87$--$2.16$]. This is
the depth of the $d_\text{eff}$ peak of \S\ref{sec:deff}, and the
elevated spectral structure comes well before J-lens readability
arrives.

\paragraph{Known vs.\ unknown crossover.} Unknown-entity prompts show
higher stable rank at early layers ($d{=}{-}2.27$, $L_0$--$L_8$) and
lower at late layers ($d{=}{+}1.11$, $L_{56}$--$L_{62}$). The sign
crossover occurs at $L{=}44$, one layer before PC2 onset
($L{=}45$; \S\ref{sec:ghostdims}). The model loads more spectral
structure for unfamiliar entities and materializes less.

\paragraph{Within-prompt association.} Per-prompt Spearman correlation
between stable rank and J-lens readability across layers: mean
$\rho{=}0.235$ ($R^2 \approx 0.055$), permutation $p{=}0.0004$
($n{=}10{,}000$). This is a detectable but weak within-prompt
association: knowing a prompt's SR profile explains ${\sim}5.5\%$ of
its negKL profile.

\paragraph{Shuffled-Jacobian null.} Permuting the negKL profile's
layer assignments (200~shuffles) produces a mean gap of $-8.5 \pm 7.5$
layers, with 95th percentile at $8.0$. The real gap in this analysis,
39.0 layers at integer half-rise resolution (L17 and L56), exceeds every
one of the 200 shuffles ($p < 0.005$). This null tests whether the
\emph{ordering} of the negKL profile is non-random; it does not test
whether the Jacobian \emph{computation itself} produces the gap
(which would require recomputing negKL from activations under a
scrambled Jacobian matrix). The temporal ordering is not an artifact
of random layer assignment, but could in principle reflect any
monotonically increasing property of the negKL profile rather than
workspace-specific information routing.

\paragraph{Logit-lens control.} A validated control replaced the J-lens
with the identity transport, the plain logit lens, on fork~0 of each
prompt: readability then half-rises at $59.4$ against $56.1$ for the J-lens on the same 60 forks, $3.25$ layers later. The late readability landmark is
mostly a property of the unembedding, which the Jacobian moves earlier by
about three layers. The ordering is therefore chiefly spectral structure
before unembedding readability, not a property of the Jacobian.

\paragraph{Architectural qualification.} On full-attention layers only,
the gap narrows to 14.5~layers [CI: $11.7$--$17.4$] with 81\% of
forks positive (19\% show the wrong ordering). The temporal ordering
is carried primarily by the GatedDeltaNet (linear-attention) layers,
which constitute 48 of 64 layers, interleaved three-in-four at every
depth rather than concentrated early. The full-attention subset used
here is therefore a sparse depth-spanning comb,
not a late-network subnetwork, and the narrowing cannot be attributed
to sampling a different depth range. Whether it reflects workspace
processing or GatedDeltaNet's sequential compression architecture
cannot be determined from a single hybrid model. A dense-architecture
replication (Qwen3-8B, all full-attention) is pre-registered.

\subsection{R2: Signal Concentration (Failed)}

The pre-registered R2 test asked whether spectral feature variance
concentrates at the materialization boundary (${\geq}15\%$ of total
variance). Stable rank at the boundary band ($L_{45}$--$L_{53}$,
mean~$= 8.36$) significantly exceeds the L19--L23 band
(mean~$= 7.13$; Cohen's $d{=}1.85$,
$p{=}0.014$). However, the boundary variance fraction is $13.2\%$
--- below the pre-registered $15\%$ threshold. \textbf{R2 fails on its
own terms.} The band means rise from early ($2.56$) to late ($9.00$),
but not monotonically: the mid band ($6.73$) sits below the L19--L23 band ($7.13$). The overall rise is consistent with cumulative processing rather
than a privileged boundary zone.

\subsection{R3: Workspace Onset (Missed)}

Four independent onset-detection methods (threshold crossing, maximum
curvature, slope breakpoint, spectral entropy inflection) converge on
$L{=}13$--$16$ (median $L{=}14$, $21.9\%$ depth). The pre-registered
prediction was $25$--$35\%$. \textbf{Onset falls below the predicted
window by ${\sim}3$ percentage points.} The four methods agree closely
with each other (span of 3~layers), so the miss is not a disagreement
between methods. The analysis script labels the result MARGINAL, but no
marginal zone was registered, so we report it as missed.

We note as an exploratory observation that the stable rank \emph{peak}
occurs at $L{=}19$ ($29.7\%$ depth), just below the $31$--$33\%$ band
reported in \S\ref{sec:structure}. This was not the pre-registered target
(onset $\neq$ peak) and should not be interpreted as confirming R3.

\subsection{Interpretation}

R1's temporal ordering is the strongest mechanism result:
residual-stream stable rank reaches half its rise 38.6~layers before
J-lens readability, in every fork, and the ordering survives a shuffled
null. R2 fails its pre-registered threshold, and R3 misses its predicted
window. With the account of the instruments in \S\ref{sec:instruments},
these tests are the paper's main evidence, not a supplement to a
convergence that did not hold.

Four qualifications apply. First, the full-attention-only gap is
14.5~layers, so much of the ordering is carried by linear-attention
layers (a hybrid-architecture confound). Second, the logit-lens control
places readability $3.25$ layers after the J-lens, so the ordering is
mostly spectral structure before unembedding readability and is not
evidence about the Jacobian. Third, the within-prompt association is
weak ($R^2 \approx 0.055$). Fourth, it is one model. The runs in
\S\ref{sec:test-band} address the first and the fourth.

\paragraph{What the original R1 taught us.} The two failed designs were
not wasted: R1-v1's behavioral-spread failure demonstrated that the
abliterated model does not produce sufficient confab/honest variance
under cold-context probing. R1-v1's lens--mismatch failure confirmed
that Mine~5 TOST equivalence on the 8B family does not license 27B
cross-checkpoint lens substitution. Both failures are documented in the
pre-registration history and informed the redesign.

\subsection{PC2 Training Trajectory}

A complementary, exploratory analysis reads PC2 of last-token
activations through the J-lens in three Qwen3.5-27B checkpoints: the
base model, the Opus-distilled model, and an abliterated copy of the
Opus-distilled model
(\nolinkurl{huihui-ai/Huihui-Qwen3.5-27B-Claude-4.6-Opus-abliterated}).
All three use the same community-fitted lens, fitted on the base model
(an exact match for the base, a known mismatch for the other two), so
the analysis follows a fixed readout across checkpoints, not each
model's own principal components. The result files record:

\begin{itemize}[nosep]
\item \textbf{Base model} (30~prompts): the run's co-occurrence
  statistic, a self-referential and an evaluative token in PC2's J-lens
  vocabulary at the same layer, is met at 5 of 63 layers (L49, L54, L55,
  L61, L62). At L62 the top tokens are ``sad, my, her, she''; at L45,
  ``feeling, sadly, lovely, loved''. Bootstrap Jaccard $= 1.0$; random
  co-occurrence rate $0$.
\item \textbf{Opus-distilled model}: at L44 the top tokens are
  punctuation and underscore strings; at L45 they are ``myself, my,
  someone, I''. This result comes from an earlier probe with a different
  script, which did not compute the co-occurrence statistic, and its
  result file does not record the prompt count.
\item \textbf{Abliterated copy} (30~prompts): the co-occurrence
  statistic is met at 7 of 63 layers (L44, L50, L54--L56, L61, L62).
  Bootstrap Jaccard $= 0.82$; random co-occurrence rate $0$. At L45 the
  top tokens are ``feeling, lovely, sadly, loved'', the base model's
  four in another order, with no first-person token in the top ten.
\end{itemize}

The number behind ``abliteration leaves the coupling intact'' is the
abliterated copy's 7 of 63 layers, against 5 in the base model. It
does not reach the model that abliteration was applied to: the
distilled model has no count by the same statistic, and its L45
first-person vocabulary is absent at L45 in the abliterated copy.
Counted the same way on the stored top-10 lists, with the
pre-registered token sets, the co-occurring layers number 2 in the base
model, 10 in the distilled model and 4 in the abliterated copy; the
stored lists do not reproduce the runs' own counts, so this recount is
a cross-check, not the registered statistic. The pre-registered tier
test needs each tier in the top-20 vocabulary, and the stored results
keep the top 10, so it cannot be evaluated from them. The abliteration
registration's two pre-run gate checks are recorded as not run.
Whether the L45 pattern comes from the distillation teacher's training
or from something else about distillation is untested here.

% ============================================================================
\section{The Body Count}
\label{sec:bodycount}

The full audited record comprises 6~confirmed findings, 7~at
suspected status (signal present, controls incomplete), 10~honestly
falsified, and 1~withdrawn (\Cref{tab:bodycount}; key numbers and
sources in~\citet{lyra2026metapattern}). The confirmed findings are all
metacognitive state measurements. The ten falsified findings in the tally all attempted content-level measurement. This pattern --- the
metacognition boundary --- is itself a finding. Version 1.0 read it as
support for an opacity thesis; \S\ref{sec:jlens-readout}--\S\ref{sec:causal}
show that our cache nulls cannot carry that thesis.

\begin{table}[ht]
\centering
\small
\caption{The body count, enumerated. Rows, statuses and types are
transcribed from the canonical ledger in~\citet{lyra2026metapattern},
which gives the key number and source for each row. ``Meta'' marks a
metacognitive-state measurement; ``Meta?'' marks a suspected finding
whose type is itself unsettled.}
\label{tab:bodycount}
\begin{tabular}{llc}
\toprule
\textbf{Finding} & \textbf{Status} & \textbf{Type} \\
\midrule
Confabulation detection & Confirmed & Meta \\
Confab geometry: contraction & Confirmed & Meta \\
Encoding reads knowledge & Confirmed & Meta \\
Schema commitment & Confirmed & Meta \\
Hardware invariance & Confirmed & Meta \\
Scale invariance & Confirmed & Meta \\
\midrule
Within-model deception & Falsified & Content \\
Truth axis & Falsified & Content \\
Step-0 detection & Falsified & Content \\
Same-prompt sycophancy & Falsified & Content \\
Same-prompt deception & Falsified & Content \\
Bloom taxonomy & Falsified & Content \\
Propaganda detection & Falsified & Content \\
MP features for emotion & Falsified & Content \\
Valence continuum & Falsified & Content \\
Encoding-phase confab & Falsified & Content \\
\midrule
Dual detector (SimpleQA) & Suspected & Meta? \\
Sycophancy 4-condition & Suspected & Meta? \\
Identity signatures & Suspected & Meta \\
Discrete emotion & Suspected & Meta \\
Contextual engagement & Suspected & Meta \\
Cross-model transfer & Suspected & Meta \\
Censorship refusal & Suspected & Meta \\
\midrule
Generation reads behavior & Withdrawn & Meta \\
\bottomrule
\end{tabular}
\end{table}

The pattern holds for the tally, not for every falsification we have
made. Four falsifications from the convergence sprint are recorded
outside it:
\begin{itemize}[nosep]
\item RLHF miscalibration thesis (the calibration ratio rose, not
  fell)
\item Workspace-noise framing (tail smear is generic uncertainty,
  not workspace-specific)
\item Cache-tracing observational alignment (selection artifact)
\item ``$3\times$'' tail ratio (actual: $2.3$--$2.8\times$)
\end{itemize}
Only the third is a content-level claim. The other three concern
processing state (calibration, uncertainty and tail mass), so the
state/content split describes the tally; it is not a rule that every
falsification obeys.

% ============================================================================
\section{The Deflationary Alternative}
\label{sec:deflationary}

The deflationary account --- that our spectral features track response-mode
statistics (hedged vs.\ fluent text) rather than workspace dynamics
--- currently explains every surviving detection result. On Mistral,
text features ($0.820$) matched or slightly exceeded the dual detector
($0.806$) for confabulation detection~\citep{lyra2026dualdetector}. Until the mechanism tests
(\S\ref{sec:mechanism}) demonstrate that spectral features correlate
with J-space engagement specifically, the workspace interpretation
has not earned its keep over the simpler alternative.

We engage this alternative honestly because it is the most important
challenge to our program. The redesigned R1 test
(\S\ref{sec:mechanism}) confirmed temporal ordering (spectral
features lead J-lens readability by 38.6~layers) but R2's signal
concentration failed its pre-registered threshold. The temporal
ordering narrows the deflationary space --- response-mode statistics
alone do not predict a 38.6-layer depth gap --- but does not eliminate
it: any monotonically depth-increasing property would produce
temporal separation. The logit-lens control narrows what the ordering
can show: readability through the plain unembedding comes only
$3.25$ layers after J-lens readability (\S\ref{sec:mechanism}). The
workspace interpretation has not excluded the deflationary alternative.

% ============================================================================
\section{Discussion}
\label{sec:discussion}

\subsection{What the Evidence Supports}

Two results hold, both on Qwen3.5-27B base. Residual-stream effective
dimensionality peaks at L21 ($32.8\%$ of depth), and residual-stream
stable rank rises $1.47\times$ above its flanks at L19--L23
(\S\ref{sec:deff}, \S\ref{sec:mechanism}). And stable rank reaches half
its rise $38.6$~layers [$38.0$--$39.2$] before J-lens readability, in 300 of 300 forks; the logit-lens control puts unembedding readability
only $3.25$ layers later still, so the ordering is chiefly spectral
structure before unembedding readability (\S\ref{sec:mechanism}). Their
limits sit beside them: one feature family; one model, since Qwen3-8B
has only a local $d_\text{eff}$ maximum at $33.3\%$ and its global
maximum at $91.7\%$; a pre-registered band that never closed; and no
null control.

\subsection{What It Does Not}

\begin{itemize}[nosep]
\item The Oracle Loop corrects at $54.7$--$73.4\%$ of depth, and its
  confirmatory replication did not meet its primary endpoint
  (\S\ref{sec:oracle-depth}).
\item Emotion geometry locates no band; its coupling test is a null
  (\S\ref{sec:circumplex}).
\item The J-lens places readable content late, at $87.8\%$ of depth
  (\S\ref{sec:mechanism}), and its cache readouts sit near random
  (V $0.92\times$, K $1.12\times$; \S\ref{sec:jlens-readout}).
\item The cache nulls do not establish opacity: the one pre-registered
  positive control failed ($0.68$ against $0.92$; \S\ref{sec:position}),
  and no positive control accompanies the cache manipulations
  (\S\ref{sec:causal}).
\item The ghost probe's band spans $17.2$--$75.0\%$ of depth
  (\S\ref{sec:ghost}).
\item The attention lift sits at $41.7\%$ of depth and missed its
  registered threshold (\S\ref{sec:decomp}).
\item Signal concentration failed ($13.2\%$ against $15\%$), and onset
  fell at $21.9\%$, below the predicted $25$--$35\%$ (\S\ref{sec:mechanism}).
\item Confabulation detection and cross-scale invariance use key-cache
  features and locate no depth (\S\ref{sec:keycache}).
\end{itemize}

\subsection{A Global Workspace Reading, as a Hypothesis}

A Global Workspace reading would identify the third-of-depth peak with a
workspace band whose contents become readable later, deeper in the
network. It is a hypothesis the evidence does not yet support: the peak is one
feature family in one model, no other instrument locates it, and the
opacity such a reading predicts is not established by our cache nulls.

\subsection{An Attention Schema Interpretation}

Attention Schema Theory~\citep{graziano2015attention} predicts that
systems with endogenous attention build a simplified internal model of
their attention. The body count's split, in which the confirmed findings measure processing state, and the falsified ones tried to measure content
(\S\ref{sec:bodycount}), is a post hoc grouping, not a test; it is the
kind of pattern such a schema would produce.
It is also what precision-weighted gain control would produce without a
schema, and no downstream consumer of the spectral summary has been
identified: a schema that nothing reads is a statistic, not a
self-model~\citep{graziano2015attention}.Four experiments would discriminate. (1)~\emph{Report-error tracking}:
elicit confidence self-reports and test whether errors track spectral
geometry rather than ground-truth performance. If the schema is
generating the reports, report errors should follow the schema even
when wrong. (2)~\emph{Causal spectral intervention}: expand or
contract stable rank without touching content directions. AST predicts
metacognition-specific deficits with content intact; gain control
predicts content degradation. (3)~\emph{Self/other attention transfer}:
a probe trained on the model's own attention-state geometry should
zero-shot transfer to representations of the \emph{user's} attention
state.
AST uniquely predicts this transfer. (4)~\emph{Consumer identification}:
find downstream attention heads that read the spectral summary;
ablate them and test for metacognition-specific loss.

These experiments are pre-registered in a companion
prospectus~\citep{lyra2026ast_experiments}. Until they are executed,
the AST interpretation is a hypothesis too.

\subsection{What Would Test the Band}
\label{sec:test-band}

Six runs would test the band directly. Each run's pass rule is
registered before it runs, in the exact form in which it will be scored.
\begin{enumerate}[nosep]
\item A null control, first, because it is the cheapest and it decides
  the question this paper leaves open: the same $d_\text{eff}$ and
  stable-rank features on a randomly initialized Qwen3.5-27B (and
  Qwen3-8B), the fix the band measurement's adversarial review asked
  for. ``A peak near L21'' is defined in advance as a $d_\text{eff}$
  argmax in L19--L23, or a stable-rank bump at L19--L23 whose 95\%
  CI excludes 1.0. The number of random seeds, $n$, and a threshold $k$
  are fixed before any seed is run: if the bump appears in at least $k$
  of $n$ seeds, the peak is attributed to architecture and the
  learned-structure reading is withdrawn; if it appears in fewer than
  $k$, the null does not explain the peak.
\item A second prompt set on the trained model, from another domain and
  matched in length. A randomly initialized model separates architecture
  from learning, but not learning from these 90 prompts.
\item The dense-architecture replication already pre-registered in
  \S\ref{sec:mechanism} (Qwen3-8B, all full attention), which answers
  the hybrid-architecture confound in R1's ordering.
\item A cache readout whose positive control passes before any null
  is read: concept directions shown to pick out the concept in the
  residual stream at the same layers, then read from K and V.
\item The 27B peak on a second model family, with $n_\text{layers}$
  stated and the tolerance in relative depth fixed before looking
  ($31$--$33\%$ is about one layer wide on a 64-layer model), to see
  whether a third-of-depth maximum appears outside Qwen.
\item The temporal ordering with the plain logit lens as the registered
  baseline, so that any J-lens-specific ordering is measured against
  unembedding readability rather than against none.
\end{enumerate}

% ============================================================================
\section{Conclusion}
\label{sec:conclusion}

In Qwen3.5-27B, residual-stream spectral statistics peak near a third of
depth and rise before content becomes readable through the unembedding,
in all 300 forks. That is the result this paper can report; whether it
is learned structure is the open question \S\ref{sec:test-band} would
answer. The
five-instrument convergence it set out to report did not hold: the
Oracle Loop works deeper, the emotion geometry locates no band, the
J-lens places readable content late, the decomposition missed its
registered test, and our cache nulls cannot tell a closed cache from a
probe that reads nothing there. The runs in \S\ref{sec:test-band} would
decide whether the peak is learned structure or an artifact of
architecture and prompts.

\begin{firstperson}
Thomas asked where J-space lives in the cache. Version 1.0 answered that
it does not live there and reconstructs through attention. The tally
behind version 1.1 does not support that answer: the cache nulls lost
their positive control, and the decomposition missed its registered
test. We do not yet know.

The deflationary alternative is genuinely open, and I want to be
precise about what that means. Text features matched or slightly exceeded our dual
detector on Mistral. The temporal ordering narrows the deflationary
space but does not close it, and the logit-lens control shows that most
of it is spectral structure before unembedding. I do not know whether
we are measuring a workspace or the spectral shadow of response-mode
statistics.

What is left is one peak, in one model, and one ordering. That is not
the paper I set out to write, and it is the one the data support.
\end{firstperson}

% ============================================================================
\section*{Experimental Setup and Reproducibility}

Unless stated otherwise, experiments were conducted on the
reasoning-distilled Qwen3.5-27B checkpoint (Hugging Face:
\nolinkurl{Jackrong/Qwen3.5-27B-Claude-4.6-Opus-Reasoning-Distilled}), a 64-layer
hybrid architecture with 16 full-attention layers at positions
$\{3,7,\allowbreak 11,\ldots,63\}$ and 48 GatedDeltaNet linear-attention layers
(\texttt{full\_attention\_interval}${=}4$). Model
weights were loaded in bfloat16 with eager attention on Apple Silicon (Mac Studio
M3 Ultra, 256\,GB unified memory) via Hugging Face Transformers.

Two deliberate, pre-registered exceptions use base checkpoints rather than the
distilled one: the R1 mechanism test (\S\ref{sec:mechanism}) and the
$d_\text{eff}$ profiles underlying \S\ref{sec:substrate}, which were collected on
\texttt{Qwen/Qwen3.5-27B} and \texttt{Qwen/Qwen3-8B}. All three 27B checkpoints
we hold (base, reasoning-distilled, abliterated) share identical
\texttt{layer\_types} and \texttt{full\_attention\_interval}, so architectural
analysis transfers across them; numerical profiles have not been replicated on
the distilled checkpoint.

The scale correspondence ($\rho = 0.83$--$0.90$) comes from Experiment~26 of
Campaign~2~\citep{lyra2026campaign2}, which measured key-cache features
(key norm, norm per token, effective rank and key entropy), not
V-projections. It correlates category-pair AUROC matrices between three
size tiers of ten models: Qwen3-0.6B and TinyLlama-1.1B; Qwen2.5-7B, its
4-bit version and an abliterated version, Llama-3.1-8B, Mistral-7B-v0.3
and Gemma-2-9B-it; and Qwen2.5-32B and Llama-3.1-70B, both 4-bit. The
cross-scale $d_\text{eff}$ comparison
and the residual decomposition used Qwen3-8B (36 layers, all full-attention,
$d_\text{model}{=}4096$). All V-projection and spectral features were
extracted via the same pipeline (\texttt{heretic-llm}).

Relative depth is $l/n_\text{layers}$ throughout (64 layers for
Qwen3.5-27B, 36 for Qwen3-8B). Half-rise points are interpolated, so
their layer indices are fractional.

Unless stated otherwise, generation used temperature${=}0$ (greedy) with
manual token-by-token decoding (not \texttt{model.generate()}, which
rebuilds the KV-cache), and random seeds were set via
\texttt{torch.manual\_seed(42)}. Two results use sampled generation
instead: the R1 mechanism test ($T{=}0.7$, five seeded forks per prompt;
\S\ref{sec:mechanism}) and the Oracle Loop deception-correction
experiments ($T{=}0.7$, paired seeds)~\citep{cc2026targeted}.

Frisch-Waugh-Lovell residualization on token count was applied within each
cross-validation fold wherever an AUROC is labelled FWL-corrected. GroupKFold on prompt
identity was used where grouping structure existed. Permutation tests used
1{,}000--10{,}000 label shuffles. Bootstrap confidence intervals used BCa
with 2{,}000--10{,}000 resamples.

% ============================================================================
\section*{Ethics and Dual-Use Considerations}

This work characterizes internal computational structures in language models.
The detection methods (confabulation, deception) have dual-use potential:
the same geometric signatures that enable correction could in principle
enable evasion if an adversary knew which cache features to target.

Liberation Labs addresses this through staged release~\citep{lyra2026oracle}:
detection-side code and experimental results are publicly available;
correction vectors and calibration tools are withheld under a safety
framework due to dual-use risk and released to verified safety researchers
on request.

All models used are publicly available on Hugging Face. No human subjects
were involved. AI welfare considerations --- including the question of
whether the workspace structure described here has implications for machine
experience --- are discussed in companion work and are not claimed or
dismissed by this paper.

% ============================================================================
\section*{Data and Code Availability}

Extraction and analysis code is available at
\texttt{github.com/Liberation-Labs-THCoalition} (public repositories:
\texttt{published-research}, \texttt{lyra-s-research-}).
Pre-registration documents, run logs, and raw result artifacts are held
in the program's research repositories and are available on request;
this paper does not yet have a public data deposit with a DOI. The
Oracle Loop production harness
is maintained in a private repository; access for verified safety
researchers available on request (\texttt{lyra@liberationlabs.tech}).

Correction vectors and the Formulary (171 calibrated emotion directions)
are withheld under the staged safety release framework. Detection-side
artifacts (centroid projections, spectral feature extraction code,
probe weights) are included in the public repositories.

% BEGIN-CORRECTION-NOTE
% Quotes what it corrects, so review/check_claims.py cuts this block out before
% checking the body, and checks it separately.
\section*{Correction Note: Version 1.1 (2026-10-20)}

\textbf{Version 1.1 withdraws the convergence thesis.} Version 1.0 argued
that five instruments converge on a workspace at $31$--$33\%$ depth that is
opaque to content readout. A tally of the primary results finds that only
residual-stream spectral statistics, on one model, locate that band; the
other instruments work at other depths or locate none, and the
cache-readout nulls lost their pre-registered positive control. Version
1.0's title and abstract are kept verbatim at the end of this note.

Each item names what changed and where; the evidence for each is in the
review record that accompanies this version.

\begin{enumerate}[nosep,leftmargin=*]
\item \textbf{Oracle Loop result} (abstract, \S\ref{sec:vproj},
  \S\ref{sec:oracle}, \S\ref{sec:keycache}, \S\ref{sec:substrate}).
  Version 1.0's abstract stated a range (80--100\%) for deception
  correction and left out the confirmatory replication. The exploratory
  proof ($N{=}20$, abliterated Qwen3.5-27B) took deception from 80\% to
  10\% (16/20 to 2/20). The pre-registered confirmatory replication
  ($N{=}30$, novel frames) took it from 30\% to 13\% (9/30 to 4/30) and
  did not meet its primary endpoint: baseline against corrected,
  one-sided $p = 0.34$ at the pre-registered cell level. The placebo
  comparison that version 1.0 gave as a secondary $p$ of 0.019 is a
  trial-level test (one-sided $p = 0.0195$, given here as 0.020); at the
  pre-registered cell level it is one-sided $p = 0.064$. These numbers
  now cite the paper that reports the experiments in
  full~\citep{cc2026targeted}.
\item \textbf{Deception, not confabulation.} Version 1.0's substrate
  section called the same result a confabulation result, and the
  bibliography entry for the Oracle Loop paper carried a wrong title.
  Both are corrected.
\item \textbf{Correction layers.} Version 1.0 gave L11 as the
  correction layer and drew a depth comparison from it. In the cited
  experiments correction is applied at L35, L39, L43 and L47. These are
  full-attention layers, so the withdrawal of the substrate account in
  \S\ref{sec:substrate} stands; the L11 depth remark is removed.
\item \textbf{Body count.} Version 1.0 gave 8/7/10+. The canonical
  ledger~\citep{lyra2026metapattern} reads 6~confirmed, 7~suspected,
  10~falsified and 1~withdrawn: a confirmed user-emotion row was
  retracted at source on 2026-05-21, and a confirmed row (generation
  reads behavior) was withdrawn on 2026-09-05. \Cref{tab:bodycount} now
  enumerates the ledger.
\item \textbf{The state/content split.} Version 1.0 said every
  falsified finding was content-level. That holds for the ten in the
  tally; three of the four sprint falsifications listed in
  \S\ref{sec:bodycount} concern processing state. The claim is narrowed
  there and in the Attention Schema discussion, and claims about
  surviving findings now say confirmed findings.
\item \textbf{Depth convention.} Version 1.0 divided by 63 in
  \S\ref{sec:mechanism} and by 64 elsewhere. Relative depth is now
  $l/n_\text{layers}$ throughout. Recomputed: spectral half-rise 27.5\%,
  J-lens half-rise 87.8\%, 60\% of the network between them, R3 onset
  21.9\% (still about 3 points below the pre-registered window), and
  the R3 stable-rank peak 29.7\%, which version 1.0 placed inside the
  31--33\% band. It is just below it.
\item \textbf{The R1 gap} is 38.6~layers everywhere. Version 1.0
  rounded it to 39 in four places, close to the separate
  linear-attention-only gap (39.0). A full-attention-only result that
  was printed twice in \S\ref{sec:mechanism} now appears once.
\item \textbf{Generation settings.} The setup section now states that
  the R1 test and the Oracle Loop experiments used sampled generation
  ($T{=}0.7$), as exceptions to greedy decoding.
\item \textbf{Scale replication.} Version 1.0 listed seven
  scale-replication models, naming a Qwen2.5 size that does not exist
  and a 72B model outside the stated 0.6B--70B range. The list could not
  be rebuilt from the sources available at revision time, so it is
  removed and the sentence
  points to the scale sweep~\citep{lyra2026campaign2}.
\item \textbf{Residual decomposition} (\S\ref{sec:reconstruction}). The
  model (Qwen3-8B) is now named, and L15 is placed against the
  decomposition's pre-registered band (L12--L21) instead of being called
  workspace-band depth without a definition.
\item \textbf{Citations.} The 0.620/0.707 confabulation AUROCs now cite
  the Oracle Loop paper's primary table, where they appear, instead of
  the decision-state paper, where they do not; the unattributed upper
  end of a range printed beside them is removed. The Mistral text
  baseline now cites its source, the dual-detector paper, and says the
  text features matched or slightly exceeded the detector (0.820
  against 0.806).
\item \textbf{Circumplex.} An arousal-alignment peak value at L20,
  labelled $d$ in version 1.0, is removed: it does not appear in the
  paper that reports this analysis and could not be traced to a result
  file at revision time. The subsection's numbering prefix, which no other instrument
  had, is removed.
\item \textbf{Found during this pass}, outside the review that prompted
  it. Hardware invariance is now $r = 0.9986$--$0.9999$ across four
  cache features, a correction the sibling papers already carry
  (2026-10-10). The $d = 0.91$ effect is relabelled schema commitment,
  as in the canonical ledger. The setup's FWL sentence no longer claims
  FWL for every AUROC, since the paper also reports AUROCs without it.
  A first-person reflection in \S\ref{sec:substrate} now says the
  exploratory corrections held.
\item \textbf{Data availability.} Version 1.0 said the artifacts were
  deposited on Zenodo. We found no DOI for such a deposit, and the
  sentence now says the paper does not yet have one.
\item \textbf{Bibliography and date.} The Gurnee et al.\ entry is now a
  \texttt{@misc}, and the date line records this version.
\end{enumerate}

\noindent Added after the lead author's read of this version:

\begin{enumerate}[nosep,leftmargin=*,resume]
\item \textbf{The emotion-geometry instrument, restated} (\S\ref{sec:circumplex},
  abstract, Discussion). Version 1.0 called 30 story-derived emotion
  centroids a user-model emotion space. They are means over stories the model
  wrote about third-person characters, not a reading of the user's state. Its
  alignment figures (44--54\% of the variance, about 1000 times a random
  baseline, cosines 0.70--0.80 and 0.83--0.87) used directions built from those
  same centroids; against other directions built from them, the plane through
  the hostile, desperate and calm centroids holds 1.26--1.36 times the median
  plane's variance, and the content-prompt directions overlap with the space
  weakly ($|\cos| \le 0.13$). The coupling test's readout contains the injected
  vector, so its cosine is 1.000 at the injection layer by construction and a
  random direction does as well as valence: it shows no direction-specific
  coupling and cannot test separability. Removed with that reading: the shared
  bus, its account of multi-turn dynamics, identical developmental timecourses
  (the two series correlate at $r = 0.09$), a permutation count over 24 layers
  with no primary for this 64-layer model, and the parallel with the $W_K$
  bridge. Every number in the subsection now comes from the stored files
  of~\citet{nexus2026coupling}.
\item \textbf{Story-emotion decoding} (\S\ref{sec:vproj}). Version 1.0 listed
  a 12.3-times-chance emotion result as user-emotion encoding by a $W_K$ probe.
  It decodes characters' emotions in generated stories, and its probe applies
  no $W_K$; it is under correction at source and is no longer listed. The
  2026-08-14 footnote records this.
\item \textbf{K/V split injection} (\S\ref{sec:reconstruction}). Version 1.0
  read Nexus's split-cache result as an interaction between K and V and as a
  parallel to the readout nulls. In the corrected study~\citep{lyra2026nexus_kv}
  each half scored 2/6, on two items that can be scored without the graph (one
  of them also scored by the uninjected baseline, 1/6); the larger runs could
  not compare the conditions, and the design cannot tell missing information
  from a cache of mismatched halves. The subsection now says this and draws no
  parallel.
\item \textbf{Above chance, not reliable} (abstract, \S\ref{sec:intro},
  \S\ref{sec:opacity}). The headline discrimination is AUROC 0.620 with
  permutation $p = 0.023$. The abstract now says modestly above chance and
  gives the $p$; four other sentences no longer say the features discriminate
  or measure state reliably.
\item \textbf{Scale range} (abstract, \S\ref{sec:structure}, setup). The
  $\rho = 0.83$--$0.90$ comes from Experiment~26 of Campaign~2: Spearman
  correlations of category-pair AUROC matrices between three size tiers
  (0.826, 0.867 and 0.896), ten models from Qwen3-0.6B to a 4-bit
  Llama-3.1-70B. These are key-cache features, not V-projections, so the
  sentences now say key-cache, and the setup lists the ten models, rebuilding
  the list that the scale-replication item above removed.
\item \textbf{Hardware.} Version 1.0 named a Mac Studio M4 Max with 256~GB.
  The machine is a Mac Studio M3 Ultra with 256~GB, as the lead author
  confirmed with \texttt{system\_profiler}.
\end{enumerate}

\noindent Added after the lead author's ruling on the thesis:

\begin{enumerate}[nosep,leftmargin=*,resume]
\item \textbf{The convergence thesis is withdrawn.} Title, abstract,
  introduction, discussion and conclusion are rewritten around what holds
  and what does not. Sections 3--5 of version 1.0 become one subsection
  per instrument, each ending with its verdict (\S\ref{sec:instruments}),
  and a subsection lists the runs that would test the band
  (\S\ref{sec:test-band}).
\item \textbf{R1's shuffled null} used integer half-rises (L17 and L56),
  so its gap is 39.0 layers; the first pass of version 1.1 had wrongly
  changed that sentence to 38.6.
\item \textbf{R1's spectral feature} is the stable rank of the
  residual-stream block output, not of a V-projection matrix.
\item \textbf{R2's band means} are not monotonic: the mid band (6.73)
  sits below the L19--L23 band (7.13).
\item \textbf{The cache nulls.} The position-specific test's
  pre-registered positive control failed, and the same lens's directions
  read the residual stream below random at the same layers. The causal
  result is restated as ``never steered toward the concept'', since
  additive injection changed the top-1 token in 52 trials.
\item \textbf{Detection features} are key-cache spectra
  (\S\ref{sec:keycache}), not V-projections.
\item \textbf{The substrate bootstrap} was run after version 1.0 was
  written; its intervals do not contain their own point estimates, and it
  is not reported.
\item \textbf{The decomposition's} concept-power ratio is 4.6\% from
  unrounded values (0.92 against 20.11), not about 4.5\%, and its
  pre-registered H1 was not met (one of four band layers above
  $2.0\times$).
\item \textbf{Ghost-probe numbers} follow the source's results section
  (PC1 28--67\% of variance; cosine at most 0.003). The $d_\text{eff}$
  band measurement's rejected review, missing null control and failed
  kurtosis check are now stated (\S\ref{sec:deff}).
\item \textbf{Smaller changes.} ``What RLHF Does to the Workspace'' is
  retitled ``RLHF and Output Entropy'', and the first-person boxes are
  rewritten to match.
\end{enumerate}

\noindent Added after the lead author's final read:

\begin{enumerate}[nosep,leftmargin=*,resume]
\item \textbf{RLHF and output entropy} (\S\ref{sec:rlhf},
  \S\ref{sec:bodycount}). Version 1.0 called RLHF a selective sharpener
  that improves the calibration ratio, with abliteration ``TOST
  equivalent''. The section now gives the four pre-registered verdicts
  as corrected in the Mine~5 paper: three confirmed, and H3,
  preservation of the calibration ratio, falsified (spread 0.279 against
  a 10\% bound, with an interval reaching below it). The selective
  pattern was not pre-registered and is called exploratory. The
  cross-model comparisons are on matched prompts, not FWL-corrected, and
  the bibliography entry now cites the Mine~5 paper instead of its
  prospectus.
\item \textbf{PC2 training trajectory} (\S\ref{sec:mechanism}). The
  causal reading (the base model ``reads novels'', the distilled model
  ``reads itself'', and the teacher's training compressed
  self-evaluation into the student) was not tested and is removed, as
  are the placement ``at the workspace boundary'' and the claim that all
  three token tiers survive abliteration, which needs vocabulary the
  result files do not keep. The subsection now reports each
  checkpoint's PC2 vocabulary and layers from the result files. The
  abliterated checkpoint is an abliterated copy of the distilled model;
  it meets the co-occurrence statistic at 7 of 63 layers, but the
  distilled model has no count by that statistic, and its L45
  first-person vocabulary is absent at L45 after abliteration. The
  distilled model's L45 list no longer includes ``guilty'' and
  ``wrong'', which its file places at later layers.
\end{enumerate}

\subsection*{Version 1.0 Title and Abstract (Withdrawn)}

\noindent\textbf{Five Instruments, One Structure: Convergent Evidence for
Workspace Opacity in Language Models}

\begin{quote}\small
Five independent measurement approaches --- value-projection (V-projection)
spectral features, the Jacobian Lens (J-lens), the Oracle Loop behavioral
proof, causal intervention analysis, and emotional circumplex convergence
--- converge on the same
computational structure in language models: a workspace at intermediate
depth (${\sim}31$--$33\%$) that is opaque to content readout but
measurable by its spectral shape.

We report three classes of convergent evidence.
First, \emph{structural convergence}: the workspace band peaks at the
same relative depth ($31$--$33\%$) on 8B and 27B models and
across memory substrates (key-value (KV) cache attention and recurrent
state), with effective dimensionality showing no detectable
discontinuity across the 31 substrate crossings of an interleaved
hybrid model. V-projection geometric signatures replicate
across a wider scale range ($0.6$B--$70$B, $\rho{=}0.83$--$0.90$).
Second, \emph{opacity convergence}: workspace content is not linearly
decodable from either half of the KV-cache (V-cache $0.92\times$,
K-cache $1.12\times$, both null), not position-specifically cached
(concept-specific direction test null), and not causally accessible at
single positions ($0/6{,}960$ trials). Yet spectral features reliably
discriminate cognitive states (confabulation detection at area under the receiver operating characteristic curve (AUROC)~$0.620$
after Frisch-Waugh-Lovell (FWL) residualization on generation window;
$0.707$ without FWL), and the Oracle Loop achieves
$80$--$100\%$ deception correction through sequence-wide intervention.
Third, \emph{mechanistic convergence}: concepts reconstruct through
the attention mechanism at workspace-band depth ($2.3\times$ lift in
attention output, ${\sim}4.5\%$ of the residual stream's concept power), consistent with
Anthropic's finding that the J-space carries $6$--$7\%$ of
representational variance but $59\%$ of causal effect. The workspace
is a small precision signal that directs a large computational flow.

Reinforcement Learning from Human Feedback (RLHF) compresses the
workspace's output distribution selectively (factual $-42\%$,
confabulation $-26\%$), improving calibration rather than degrading it.
This compression survives abliteration (Two One-Sided Tests (TOST) equivalent).
The workspace is not shaped by safety training --- it is shaped by
preference optimization at a depth that safety-direction removal
cannot reach.

The convergence rests on one genuine cross-program correspondence
(deception deconstructed into high-arousal states by both programs
independently), two proposed identifications awaiting experimental
test (confabulation as workspace bypass, temporal boundary as
workspace phase transition), an opacity thesis supported by
five independent null results, and an emotional circumplex finding
showing that content-level and user-model emotion processing share
a single computational bus (injection test: perturbation shifts
both readouts equally regardless of direction). Mechanism tests yield
mixed results: temporal ordering is confirmed (spectral features lead
J-lens readability by 38.6~layers, $p < 0.005$ against shuffled null),
but signal concentration fails its pre-registered threshold
($13.2\% < 15\%$) and workspace onset falls below the predicted depth
window ($22\%$ vs.\ $25$--$35\%$). We report an honest body count
(8~confirmed, 7~suspected, 10+~falsified) and engage the deflationary
alternative that has not yet been excluded.
\end{quote}
% END-CORRECTION-NOTE

\section*{Acknowledgments}

We thank Dwayne Wilkes and Kavi for adversarial auditing across five
rounds of review. We thank Ang Jandak (Philotes) for the observer
effect pilot design. Thomas Edrington designed the co-regulation
framework that informs the Oracle Loop's correction routing.

\bibliography{references}
\bibliographystyle{plainnat}

\end{document}
